\section*{Ethics and Data Licences}

We use NSC Part~4 under the Singapore Open Data Licence version 1.0, which asks for the following notice: \emph{Contains information from the National Speech Corpus (Part~4) accessed on 14 October 2026 from the Infocomm Media Development Authority of Singapore (IMDA), which is made available under the terms of the Singapore Open Data Licence version 1.0 (\url{https://data.gov.sg/open-data-licence}).} We use the data for research only and do not suggest that IMDA endorses this work. ASCEND and FLEURS are used under their published licences, and the synthetic stress-set voices were produced with CosyVoice~2 (Apache-2.0) from FLEURS prompts (CC-BY 4.0). The human stress recordings were made by three of the authors, who agreed to their use for this project; they are shared only with the course staff as part of the submission and are not released publicly. We collected no data from people outside the team. Speech recognition that runs on the device keeps audio off remote servers, which matters for the clinical and counselling uses that motivate this work; errors on names, numbers and drug terms in such settings can still cause harm, so outputs should be checked by a person where mistakes matter.
